% =====================================================================
%  CHAPTER 5 — CONCLUSION AND FUTURE WORK
%  Graduation thesis: A hybrid (signature + AI) web intrusion detection
%  system operating on Apache HTTP/HTTPS access logs.
%
%  HOW TO USE
%  ----------
%  * Standalone build: compiles on its own (run pdflatex twice for refs).
%    As a standalone document the chapter is numbered 1; inside the thesis
%    it becomes the final chapter automatically.
%  * Integration: delete the two "STANDALONE WRAPPER" blocks and keep the
%    content from "\chapter{...}" onward. The \cite{...} keys resolve once
%    this chapter is included alongside the others and the shared
%    bibliography is built.
% =====================================================================

% --------------------- STANDALONE WRAPPER (begin) --------------------
\documentclass[12pt,a4paper]{report}
\usepackage[utf8]{inputenc}
\usepackage[T1]{fontenc}
\usepackage{lmodern}
\usepackage{amsmath}
\usepackage{enumitem}
\usepackage{url}
\usepackage[hidelinks]{hyperref}
\usepackage{geometry}
\geometry{margin=1in}
\usepackage{textcomp}

\newlist{romanlist}{enumerate}{1}
\setlist[romanlist]{label=(\roman*), leftmargin=2.2em, itemsep=2pt}

% Note: the \cite keys (pham2016, althubiti2017) resolve from the global
% bibliography once this chapter is integrated. In a standalone build they
% will show as undefined-citation placeholders, which is harmless.

\begin{document}
% ---------------------- STANDALONE WRAPPER (end) ---------------------

\chapter{Conclusion and Future Work}
\label{chap:conclusion}

This final chapter closes the thesis. Section~\ref{sec:summary} summarises
the problems the project set out to solve and what it achieved, and then
states honestly the limitations that remain. Section~\ref{sec:future} builds
on those limitations to propose concrete directions for future development.

% =====================================================================
\section{Summary}
\label{sec:summary}

The objective of this project was to build a web intrusion detection system
that overcomes the central weakness of the existing literature surveyed in
Chapter 2, namely the reliance on a single detection
paradigm. Signature-based detectors are precise but blind to new attacks,
purely supervised detectors require labelled data and cannot generalise
beyond the attack families they were trained on, and purely unsupervised
detectors catch novel anomalies but at the price of many false alarms. The
project addressed this gap by designing, implementing, and evaluating a
three-tier hybrid system that fuses all three paradigms, and the work
delivered the following concrete results.

\begin{romanlist}
  \item \textbf{A working three-tier hybrid detector.} The system combines a
        signature- and rule-based tier, a supervised Random Forest tier, and
        an unsupervised one-class anomaly tier, and reconciles their opinions
        through a smart consensus rule in which the supervised model
        validates the anomaly tier. This design, detailed in
        Chapter 3, is the core contribution of the thesis.
  \item \textbf{A balanced benchmark dataset.} A corpus of $111{,}065$ log
        lines was assembled from the CSIC 2010 dataset, real attack payloads
        from a public repository, and synthetic traffic, balanced evenly
        between attack and clean requests, and partitioned so that the
        one-class tier could be trained on clean traffic alone.
  \item \textbf{Operation on raw server logs.} Unlike most prior work, which
        operates on pre-processed feature sets, the system works directly on
        Apache access logs. It performs multi-layer URL decoding to defeat
        encoding-based evasion, supports attacks delivered through the POST
        body, and represents every request with a compact, shared
        $22$-feature vector.
  \item \textbf{Strong, recall-oriented results.} As reported in
        Chapter 4, the supervised tier alone reaches an F1
        of $94.78\%$ and the one-class tier alone an F1 of $91.57\%$, while
        the proposed smart consensus achieves the system's best F2 score of
        $96.26\%$, combining a recall of $97.62\%$ with a precision of
        $91.19\%$. The consensus therefore meets the operational priority of
        an intrusion detection system, which is to minimise missed attacks.
  \item \textbf{Explainability and a usable tool.} Because the verdict is a
        transparent combination of three named signals, every alert is
        annotated with the tiers that produced it, and the system ships with
        a real-time monitoring mode and a dashboard, making it usable in
        practice rather than only as an offline experiment.
\end{romanlist}

Despite these results, several limitations remain and should be acknowledged.
First, the models were trained on the CSIC 2010 e-commerce traffic together
with generic web traffic, so when deployed on an application with very
different normal behaviour (for example a content-management system or a
bespoke API) the system may raise more false positives until it is retrained
on that application's own clean logs. Second, the unsupervised tier assumes a
stable notion of ``normal'' traffic; if the legitimate behaviour of the
protected application drifts over time, previously unseen but benign patterns
may be flagged as anomalies, a problem only partially mitigated by the smart
consensus. Third, the anti-evasion decoding handles a bounded number of
encoding layers, so payloads hidden behind an unusually deep or binary
encoding can still slip through. Fourth, the system inspects textual request
content, so attacks carried in binary POST bodies are not analysed. Finally,
the evaluation, although conducted on a sizeable balanced set, was performed
on a single dataset family and on logs already decrypted by the server; the
system performs detection and alerting only and does not, by itself, block
malicious requests. These limitations directly motivate the directions
proposed next.

% =====================================================================
\section{Suggestion for Future Works}
\label{sec:future}

Building on the limitations identified above, several avenues for future
development are worth pursuing.

\begin{romanlist}
  \item \textbf{Evaluation on more and more realistic datasets.} The most
        important next step, also recommended by Pham et al.~\cite{pham2016},
        is to validate the system on additional datasets, especially traffic
        collected from a real production environment, so that its
        generalisation beyond CSIC 2010 can be measured rather than assumed.

  \item \textbf{Semi-supervised and online learning.} Following the
        suggestion of Althubiti et al.~\cite{althubiti2017}, semi-supervised
        techniques could reduce the labelling cost while retaining accuracy.
        In parallel, an online or incremental learning scheme that
        periodically refreshes the models on recent clean traffic would
        address the baseline-drift limitation of the unsupervised tier.

  \item \textbf{Concept-drift detection and an analyst feedback loop.} The
        system could monitor its own false-positive rate and trigger
        retraining when drift is detected, and it could incorporate the
        analyst's confirmations and dismissals of alerts as new labels,
        turning day-to-day operation into a continuous improvement process.

  \item \textbf{Richer representations for the anomaly tier.} Replacing or
        augmenting the hand-crafted feature vector with a learned
        representation, for instance from an autoencoder or a sequence model
        over the request text, could let the unsupervised tier capture
        semantic anomalies that the current features miss, while preserving
        the one-class training discipline.

  \item \textbf{Stronger evasion and adversarial resistance.} Extending the
        decoder to handle deeper and binary encodings, and explicitly testing
        the system against adversarially crafted payloads, would harden it
        against attackers who deliberately target the detector.

  \item \textbf{From detection to prevention.} The system could be
        integrated with a web application firewall so that high-confidence
        alerts trigger automatic blocking, turning the present detection-only
        pipeline into an active prevention mechanism, with the consensus
        score serving as a natural, tunable threshold for the blocking
        decision.

  \item \textbf{Production deployment at scale.} In a real deployment the
        detector would operate inside a streaming pipeline rather than
        process a static log file, and it would be combined with several
        standard infrastructure components. Web servers (Apache or Nginx)
        would ship their access logs through a lightweight forwarder such as
        Filebeat or Fluent Bit into an \emph{Apache Kafka} topic that acts as
        a buffering queue; this decouples the rate at which logs are produced
        from the rate at which they are analysed and absorbs sudden traffic
        spikes without losing events. The IDS itself would be packaged as a
        container and run as one or more Kafka consumers on a
        \emph{Kubernetes} cluster: when Kafka reports that the queue is
        growing (a traffic surge), a horizontal autoscaler---for example KEDA
        driven by the Kafka consumer lag---would automatically replicate the
        IDS from a single container to many parallel ones, and scale them
        back down once the backlog drains, so the system neither loses data
        nor wastes resources. The JSON alerts emitted by the consensus stage
        would then be forwarded to a \emph{SIEM} platform (such as Elastic/ELK,
        Splunk, Wazuh, or IBM QRadar) for correlation with other security
        events, long-term storage, dashboards, and incident response; the
        highest-confidence alerts could additionally drive the web application
        firewall of the previous point to block offending requests in real
        time. Designing, implementing, and benchmarking this architecture
        under production load is a natural direction for future engineering
        work.
\end{romanlist}

Taken together, these directions would extend the thesis from a strong,
explainable, recall-oriented detector validated on a benchmark dataset
towards an adaptive, self-improving, and preventive system suitable for
long-term deployment in a changing production environment.

% --------------------- STANDALONE WRAPPER (begin) --------------------
\end{document}
% ---------------------- STANDALONE WRAPPER (end) ---------------------
